\SourcePage{632}{635}
\section{Analytic properties of the scattering amplitude}

Several important properties of the scattering amplitude can be established
by studying it as a function of the scattered particle's energy $E$, formally
regarded as a complex variable.

Consider motion in a field $U(r)$ which tends sufficiently rapidly to zero at
infinity; the required rate will be specified below. To simplify the argument,
suppose initially that the particle's orbital angular momentum is $l=0$. Write
the asymptotic form

\SourcePage{633}{636}
of the wave function, a solution of the Schrödinger equation with $l=0$ at an
arbitrary specified $E$, as
\begin{equation}
 \chi=r\psi=A(E)\exp\left(-\frac{\sqrt{-2mE}}{\hbar}r\right)
 +B(E)\exp\left(\frac{\sqrt{-2mE}}{\hbar}r\right),
 \tag{128.1}
\end{equation}
and regard $E$ as complex. We define $\sqrt{-E}$ to be positive when $E$ is
real and negative. The wave function is assumed normalized by some fixed
condition, say $\psi(0)=1$.

On the left-hand part of the real axis ($E<0$), the exponential factors in the
two terms of (128.1) are real; one decreases and the other increases as
$r\to\infty$. Since $\chi$ is real, $A(E)$ and $B(E)$ are real for $E<0$.
Consequently their values at any two points symmetric about the real axis are
complex conjugates:
\begin{equation}
 A(E^*)=A^*(E),\qquad B(E^*)=B^*(E).
 \tag{128.2}
\end{equation}

Continuing from the negative to the positive real semiaxis through the upper
half-plane gives, for $E>0$,
\begin{equation}
 \chi=A(E)e^{ikr}+B(E)e^{-ikr},\qquad k=\frac{\sqrt{2mE}}{\hbar}.
 \tag{128.3}
\end{equation}
Continuation through the lower half-plane would instead give
\[
 \chi=A^*(E)e^{-ikr}+B^*(E)e^{ikr}.
\]
Since $\chi$ must be single-valued in $E$, it follows that
\begin{equation}
 A(E)=B^*(E)\quad\text{for }E>0.
 \tag{128.4}
\end{equation}
This also follows directly from the reality of $\chi$ for $E>0$. The
coefficients $A(E)$ and $B(E)$ themselves are nevertheless multivalued because
of the square root $\sqrt{-E}$ in (128.1). Remove this ambiguity by cutting the
complex plane along the positive real semiaxis. The cut makes $\sqrt{-E}$, and
hence $A(E)$ and $B(E)$, single-valued. Their values on the upper and lower
edges are complex conjugates; in (128.3), $A$ and $B$ are evaluated on the
upper edge.

\SourcePage{634}{637}
The complex plane cut in this way will be called the \emph{physical sheet} of
the Riemann surface. With our definition,
\begin{equation}
 \mathop{\rm Re}\sqrt{-E}>0
 \tag{128.5}
\end{equation}
throughout this sheet. In particular, on the upper edge of the cut,
$\sqrt{-E}$ becomes $-i\sqrt E$.

In (128.3), the factors $e^{ikr}$ and $e^{-ikr}$, and therefore the two terms
in $\chi$, have the same order of magnitude, so this asymptotic expression is
always valid. Everywhere else on the physical sheet, the first term in (128.1)
decreases exponentially and the second increases as $r\to\infty$, by (128.5).
The terms then have different orders, and (128.1) may cease to be a legitimate
asymptotic form: the smaller term against the larger may exceed the permitted
accuracy. For (128.1) to be valid, the ratio of the small term to the large one
must not be below the relative order $U/E$ neglected in the Schrödinger
equation on passing to the asymptotic region. Thus $U(r)$ must decrease as
$r\to\infty$ faster than
\begin{equation}
 \exp\left(-\frac{2\sqrt{2m}}{\hbar}r\mathop{\rm Re}\sqrt{-E}\right).
 \tag{128.6}
\end{equation}
If this holds for every $\mathop{\rm Re}\sqrt{-E}>0$, that is, if $U(r)$
decreases faster than
\begin{equation}
 e^{-cr}
 \tag{128.6a}
\end{equation}
for every positive constant $c$, an asymptotic expression of the form (128.1)
is valid throughout the physical sheet. As a solution of an equation with
finite coefficients, it has no singularities in $E$. Hence $A(E)$ and $B(E)$
are regular throughout the physical sheet except at $E=0$, the end of the cut
and therefore their branch point.\footnote{Throughout this section we study
the scattering amplitude on the physical sheet. Later, however, it will
sometimes be necessary to consider the second, \emph{unphysical} sheet of the
Riemann surface (see \S~134). On that sheet
\begin{equation}
 \mathop{\rm Re}\sqrt{-E}<0.
 \tag{128.5a}
\end{equation}
The unphysical sheet is reached directly downward through the cut from the
positive semiaxis.}

\SourcePage{635}{638}
Bound states in $U(r)$ have wave functions tending to zero as $r\to\infty$.
The second term of (128.1) must therefore vanish, so the discrete energy levels
are zeros of $B(E)$. Since the Schrödinger equation has only real eigenvalues,
all zeros of $B(E)$ on the physical sheet are real and lie on the negative real
semiaxis.

For $E>0$, $A(E)$ and $B(E)$ are directly related to the scattering amplitude.
Comparing (128.3) with (33.20), written as
\begin{equation}
 \chi=\mathrm{const}\cdot[e^{i(kr+\delta_0)}-e^{-i(kr+\delta_0)}],
 \tag{128.7}
\end{equation}
gives
\begin{equation}
 -\frac{A(E)}{B(E)}=e^{2i\delta_0(E)}.
 \tag{128.8}
\end{equation}
By (123.15), the $l=0$ scattering amplitude is
\begin{equation}
 \hbar f_0=\frac1{2ik}(e^{2i\delta_0}-1)
 =\frac{\hbar}{2\sqrt{-2mE}}\left(\frac AB+1\right),
 \tag{128.9}
\end{equation}
where $A$ and $B$ are taken on the upper edge of the cut.

Considering the amplitude as a function of $E$ over the physical sheet, we
see that the discrete energy levels are its simple poles. If $U(r)$ satisfies
(128.6a), the amplitude has no other singular points, by the preceding
argument\footnote{Apart from $E=0$, which is singular because of the stated
singularity of $A(E)$ and $B(E)$. The scattering amplitude itself remains
finite as $E\to0$ (see \S~132). For brevity, this qualification will not be
repeated below.}.

Let us calculate the residue at a pole corresponding to a discrete level
$E=E_0<0$. The equations for $\chi$ and its energy derivative are
\[
 \chi''+\frac{2m}{\hbar^2}(E-U)\chi=0,\qquad
 \left(\frac{\partial\chi}{\partial E}\right)''
 +\frac{2m}{\hbar^2}(E-U)\frac{\partial\chi}{\partial E}
 =-\frac{2m}{\hbar^2}\chi.
\]
Multiply the first by $\partial\chi/\partial E$, the second by $\chi$,
subtract, and integrate over $dr$. One obtains
\begin{equation}
 \chi'\frac{\partial\chi}{\partial E}
 -\chi\left(\frac{\partial\chi}{\partial E}\right)'
 =\frac{2m}{\hbar^2}\int_0^r\chi^2dr.
 \tag{128.10}
\end{equation}

\SourcePage{636}{639}
Apply this at $E=E_0$ and $r\to\infty$. The integral on the right becomes
unity if the bound-state wave function is normalized by
$\int\chi^2dr=1$. In the left-hand side insert (128.1), noting that near
$E=E_0$
\[
 A(E)\approx A(E_0)=A_0,\qquad
 B(E)\approx(E+|E_0|)\left.\frac{dB}{dE}\right|_{E=E_0}
 =\beta(E+|E_0|).
\]
The result is $\beta=-1/(A_0\hbar\sqrt{2m|E_0|})$.

These expressions show that near $E=E_0$ the leading term of the scattering
amplitude, which is the $l=0$ amplitude, is
\begin{equation}
 f=\frac{\hbar^2A_0^2}{2m}\frac1{E+|E_0|}.
 \tag{128.11}
\end{equation}
Thus its residue at a discrete level is determined by the coefficient $A_0$
in the asymptotic expression
\begin{equation}
 \chi=A_0\exp\left(-\frac{\sqrt{2m|E_0|}}{\hbar}r\right)
 \tag{128.12}
\end{equation}
for the normalized wave function of the corresponding stationary state.

Return now to the analytic properties when (128.6a) is not fulfilled. In such
fields only the increasing term in (128.1) is a valid part of the asymptotic
solution throughout the physical sheet. Accordingly, $B(E)$ can still be said
to have no singularities.

Under these conditions $A(E)$ can be defined in the complex plane only by
analytic continuation of the coefficient in the asymptotic expression for
$\chi$ on the positive real semiaxis, where both terms are valid. Such a
continuation generally gives different results according as it is made from
the upper or lower edge of the cut.

To make the definition single-valued, define $A(E)$ in the upper and lower
half-planes by analytic continuation from the upper and lower sides,
respectively, of the positive semiaxis. In general the cut must then

\SourcePage{637}{640}
be extended along the whole real axis. The function so defined still satisfies
$A(E^*)=A^*(E)$, but need not be real on either the positive or negative real
semiaxis. In principle it may also have singularities.

Nevertheless, there is a class of fields for which $A(E)$ has no singularities
inside the physical sheet even though (128.6a) does not hold.

To see this, regard $\chi$ as a function of complex $r$ at a fixed complex
$E$. It is enough to consider $E$ in the upper half-plane, since the values of
$A(E)$ in the two half-planes are complex conjugates. For values of $r$ such
that $Er^2$ is real and positive, the two terms in (128.1) have the same order.
This restores the situation for positive $E$ and real $r$, in which both
asymptotic terms are valid for any field tending to zero at infinity. Thus
$A(E)$ cannot have singular points at those $E$ for which $U(r)\to0$ as
$r\to\infty$ along the ray on which $Er^2>0$. As $E$ ranges over the upper
half-plane, this condition selects the lower right quadrant of the complex
$r$-plane. We conclude that $A(E)$ has no singularities within the physical
sheet also when\footnote{Since $U(r)$ is real on the real axis,
$U(r^*)=U^*(r)$. Thus fulfillment of (128.13) in the lower right quadrant
automatically implies its fulfillment throughout the right half-plane.}
\begin{equation}
 U(r)\to0,\quad\text{as }r\to\infty\text{ in the right half-plane}
 \tag{128.13}
\end{equation}
(L.~D.~Landau, 1961).

Conditions (128.6a) and (128.13) cover a very broad class of fields. As a rule,
therefore, the scattering amplitude has no singularities inside either
half-plane. On the negative semiaxis itself, which belongs to the physical
sheet when there is no cut there, it has poles at the bound-state energies; if
a cut is present, other singularities may also occur there.

\SourcePage{638}{641}
The latter situation occurs, in particular, for fields of the form
\begin{equation}
 U=\mathrm{const}\cdot r^n e^{-r/a}
 \tag{128.14}
\end{equation}
for any $n$. On the segment $0<-E<\hbar^2/(8ma^2)$ of the negative semiaxis,
(128.6) holds; there can be no cut there, and the amplitude has only the poles
of bound states. On the remainder of the negative semiaxis there may also be
\emph{extraneous} poles and other singularities (S.~T.~Ma, 1946). They arise
because (128.14) ceases to tend to zero along the ray $Er^2>0$ as soon as $E$
passes below the negative semiaxis, when that ray passes leftward beyond the
imaginary axis in the complex $r$-plane.

Next consider the analytic properties as $|E|\to\infty$. When
$E\to+\infty$ along the real axis, the Born approximation applies and the
scattering amplitude tends to zero. The same situation holds when $E$ tends
to infinity in the complex plane along a line $\arg E=\mathrm{const}$, if one
considers complex $r$ for which $Er^2>0$. If $U\to0$ along
$\arg r=-(1/2)\arg E$ and has no singular points on that line, the Born
approximation is applicable and the amplitude again tends to zero. As
$\arg E$ ranges from $0$ to $\pi$, $\arg r$ ranges from $0$ to $-\pi/2$.

It follows that the scattering amplitude tends to zero at infinity in every
direction of the $E$ plane if $U(r)$ has no singular points in the right
half-plane of $r$ and tends to zero at infinity there.

Although the discussion was phrased for $l=0$, all these results also apply to
partial scattering amplitudes with any nonzero angular momentum. The only
difference in the derivation is that the exact free radial wave functions
(33.16) must replace the factors $e^{\pm ikr}$ in the asymptotic expressions
for $\chi$\footnote{The limiting form (33.17) of those functions may be used
only for $E>0$. Elsewhere in the $E$ plane the two terms in $\chi$ have
different orders; using the limiting expressions would generally introduce an
error larger than that due to neglecting $U$ in the Schrödinger equation.}.

\SourcePage{639}{642}
For $l\ne0$, formulas (128.9) and (128.11) require some modifications. Instead
of (128.7), one has
\begin{equation}
 \chi_l=rR_l=\mathrm{const}\cdot\left\{
 \exp\left[i\left(kr-\frac{l\pi}2+\delta_l\right)\right]
 -\exp\left[-i\left(kr-\frac{l\pi}2+\delta_l\right)\right]\right\}
 \tag{128.15}
\end{equation}
and the partial amplitude $f_l$, defined as in (123.15), is
\begin{equation}
 f_l=\frac{\hbar}{2\sqrt{-2mE}}
 \left[(-1)^l\frac AB+1\right].
 \tag{128.16}
\end{equation}
In place of (128.11), the leading term of the scattering amplitude near a
level $E=E_0$ of angular momentum $l$ is
\begin{equation}
 f\approx(2l+1)f_lP_l(\cos\theta)
 =(-1)^{l+1}\frac{\hbar^2A_0^2}{2m}
 \frac1{E+|E_0|}(2l+1)P_l(\cos\theta).
 \tag{128.17}
\end{equation}
